In this chapter, we will study the most fundamental unit of building SNNs -- the \textbf{Spiking Neurons}. When it comes to the SNNs, you may think of \textsl{Spiking Neurons} as \textit{analogous} to the \textsl{Artificial Neurons} in the DNNs. This chapter starts with a section on the \textsl{Biological Neurons} that lay down the neuroscience principles underlying the spiking neurons. Later, a section on the \textit{functional comparison} between spiking neurons and artificial neurons is also presented, where we highlight the inherent temporality and sparsity of spiking neurons.

\subsection{Should we write a section on why Spiking Neurons}
Sparse method of encoding and working with temporal information. H/W neurons don't need to be active all the time.

\section{Dales Principle}
and How SNNs violate that?

\section{Biological Neurons} \label{sec:bio-neuron}
\subsection{Action Potentials}
Perhaps look here: 

\url{https://pure.rug.nl/ws/portalfiles/portal/1106997626/Complete_thesis.pdf}

Introduce the notations for membrane potential, V[t], input current I[t], etc. here and then explain them in the Spiking Neurons section.


\newpage
\section{Spiking Neurons}
\textbf{Spiking Neurons} are electro-mathematical abstraction of biological neurons. They are generally \textit{not} very detailed representations of biological neurons, rather simple enough to reproduce their intended spiking behaviour. Note that a majority of the spiking neurons you will come across will fall into the category of \textsl{point} spiking neuron models; few rare ones will fall into the category of \textsl{spatial} spiking neuron models. The \ul{major characteristics} of biological neurons -- that are of common interest to mimic (via spiking neurons) are:
\begin{itemize}
    \item \ul{Accounting incoming action-potentials}: % https://www.perplexity.ai/search/in-the-biological-neurons-what-47xWVzh_Td.1IE1ppBqXYQ
    Spiking neurons emulate this behavior by \textit{integrating} the incoming action-potentials into their membrane potential/voltage -- either in a \textit{decaying} or \textit{non-decaying} fashion (more details later). Assuming the incoming action-potentials positively contribute, the spiking neuron's potential/voltage increases with time and eventually reaches/crosses a certain set voltage threshold.
    \vspace{5pt}
    
    \item \ul{Generating an output action-potential}: %Biological neurons generate action potential by aggregating inputs from its \textit{dendrites} on the \textit{soma}. 
    Spiking neurons emulate the generation of action-potential by %aggregating the inputs from its pre-synaptic neurons and 
    producing a \textit{binary}/\textit{graded} \textsl{spike}, where a \textit{binary} \textsl{spike} implies a binary value $\in \{0, 1\}$ and a \textit{graded} \textsl{spike} implies an integer value $\in \mathbb{Z}^+ > 1$. Note that the values of spikes are also sometimes referred as their \textit{amplitude}.
    \begin{remark}
        Some SNN implementations may use negative spikes! Also, it is commonly agreed that biological neurons' action potentials do \textit{not} have the notion of amplitude.
    \end{remark}
    \vspace{5pt}
    
    \item \ul{Resetting the membrane potential}: Spiking neurons emulate the resetting of membrane potential/voltage via two common methods: \textsl{hard-reset} and \textsl{soft-reset}, %where the membrane potential/voltage of a neuron that spiked (i.e., generated a spike) is immediately \textit{reset}. 
    where \textsl{hard-reset} implies setting the neuron's voltage to $0$, whereas \textsl{soft-reset} implies setting the neuron's voltage to a value that is equal to the neuron's current voltage \textit{subtracted} by its assumed voltage threshold. The difference between these two will be clear in the later sections.
    \vspace{5pt}
    
    \item \ul{Entering into refractory state}: Spiking neurons emulate this behavior by generally keeping their membrane potential/voltage at $0$ (in case of soft-reset) or at a subtracted value (in case of hard-reset) for a certain number of time-steps. This characteristic is typically ignored while building SNNs and is subject to varied implementations.
    \vspace{5pt}
    
    \item \ul{Propagating the action-potential along axon}: Spiking neurons generally do \textit{not} emulate this behaviour, except for the \textsl{spatial} spiking neuron models. This characteristic/property is typically embedded into the SNNs itself via the concept of \textit{delays}. In other words, when an action-potential propagates along the neuron's axon, it reaches the axon-terminal typically after some delays; this behaviour is accommodated in SNNs via introducing delays in spike propagation between the pre-synaptic and post-synaptic layers.
\end{itemize}

\subsection{Point Spiking Neuron Models}
\textbf{Point spiking neurons} are minimalistic neuron models that use simplified electro-mathematical equations to mimic the spike generation, voltage reset, and refractory state behaviour of biological neurons; they conveniently ignore to model the ionic channels, axial conductance, axonal propagation of spike/action-potential, etc., basically anything that relates to the spatial form of a biological neuron. Common examples of point spiking neurons are \textsl{Integrate \& Fire}, \textsl{Leaky Integrate \& Fire}, and \textsl{Resonate \& Fire} neuron models. 
\subsubsection{Integrate \& Fire neuron model}
\textbf{Integrate \& Fire} (IF) neuron model is the simplest of all the spiking neuron models -- by virtue of which, it's dynamic behaviour is also quite limited. An IF neuron, as the name goes, (at the very least) \textit{integrates} the incoming/input spikes into its \textit{non-decaying} membrane potential/voltage, followed by generating an output spike when its voltage reaches/crosses the set voltage threshold; whether or not it undergoes through a refractory state, is dependent on the subjective implementation. Following is the \textit{continuous-time} equation of an IF neuron: 
\begin{align}
C\frac{dV(t)}{dt} &= I(t) \\
\end{align}
Note the inherent temporality in $t$.

To implement an IF neuron on digital systems, i.e., your computer, it is necessary to discretize its continuous-time equations. Following are the \textit{discrete-time} equations (i.e., Eqs \eqref{eq:discrete-if-cur-update}, \eqref{eq:discrete-if-vol-update}, \eqref{eq:discrete-if-spk-out}, \& \eqref{eq:discrete-if-v-reset}) that describe an IF neuron: 

\begin{highlight}
\vspace{-10pt}
\begin{align}
I[t] &= (1 - \tau_\text{cur})\times I[t-1] + w\times S_\text{inp}[t] \label{eq:discrete-if-cur-update} \\
V[t] &= V[t-1] + I[t] \label{eq:discrete-if-vol-update}\\
S_\text{out}[t] &= \Theta(V[t] - V_\text{thr}) \label{eq:discrete-if-spk-out} \\
V[t] &\leftarrow V_\text{rest} \text{\qquad\qquad$\cdots$ if $V[t]>V_\text{thr}$} \label{eq:discrete-if-v-reset}
\end{align}
\end{highlight}

Before we begin explaining them, consider a \textsl{Spike Generator} that generates binary spikes $S_\text{inp}[t]$ (i.e., $S_\text{inp}[t]$ at time-step $t$ can either be $0$ or $1$) and feeds them to the IF neuron (see Fig \ref{fig:spk-neuron:if}). Note that spike generators are mere \textit{programming constructs} (that follow a desired implementation) to generate spikes and stimulate the connected neuron(s).

\begin{figure}[h]
    \centering
    \begin{tikzpicture}[
    spk-gen/.style={draw, star, very thick, fill=red!50, minimum width=0.5cm, scale=1.25},
    neuron/.style={draw, circle, very thick, fill=orange!40, scale=2, thick, inner sep=1pt}
    ]
    % Draw nodes.
    \node[spk-gen, label=below:\shortstack{\texttt{Spike}\\\text{Generator}}] (SpkGen) at (0,0) {};
    \node[neuron, label=below:\texttt{IF Neuron}] (IF-neuron) at (4,0) {\tiny $\curlywedge$};

    % Draw connections.
    \path[->] (SpkGen) edge [very thick] node[below] {$w$} (IF-neuron);
    
    \end{tikzpicture}
\caption{A \texttt{Spike Generator} stimulating an \texttt{IF Neuron}. $w$ is the weight of connection.}
\label{fig:spk-neuron:if}
\end{figure}

% \setcolumnwidth{0.7\textwidth, 0.3\textwidth}
% \begin{paracol}{2}
% First column:
% \switchcolumn
% Second column:
% \end{paracol}

In the Eq \eqref{eq:discrete-if-cur-update}, $I[t]$ and $I[t-1]$ are the IF neuron's \textit{current} at time-step $t$ and $t-1$, $\tau_\text{cur}$ is the \textit{current decay} constant (s.t., $0\leq\tau_\text{cur}\leq1$). As can be seen in the \textit{current update} Eq \eqref{eq:discrete-if-cur-update}, the current at time-step $t$ i.e., $I[t]$ is a sum of the \textit{decayed} value of current at the previous time-step $t-1$ i.e., $I[t-1]$ (decayed by the factor $(1 - \tau_\text{cur})$) \ul{and} the $w$ weighted spike input $S_\text{inp}[t]$.

In the Eq \eqref{eq:discrete-if-vol-update}, 


where, $S_\text{inp}[t]$ are the \textit{input} spikes from the preceding/pre-synaptic neuron to the IF neuron, $I[t]$ and $V[t]$ are the IF neuron's \textit{current} and \textit{voltage} respectively -- all at the \textit{discrete} time-step $t$. Note that $I[t-1]$ and $V[t-1]$ are the IF neuron's current and voltage at the previous time-step `$t-1$' -- accounting these values realizes the stateful-ness of the IF neuron. Also note that the current $I[t]$ intakes a , where , whereas, there is no such decay term for $V[t]$'s update.

\subsubsection{Leaky Integrate \& Fire Neuron Model}
\begin{remark}
    Most of the literature on SNNs use point neuron models because they are simple and easy to implement in both software and hardware.
\end{remark}

\begin{highlight}
Some highlight
\end{highlight}
\subsection{Spatial Spiking Neuron Models}
\subsubsection{Hodgkin Huxley Neuron Model}

\begin{remark}
Here is some remark \ref{sec:proin}.
\end{remark}

\newpage % Otherwise the paragraph overlaps with the text in the vertical ribbon.
\paragraph{Spiking Neurons Summary}
Example of a website \href{https://example.com/}{example.com}.